\subsection{Upper bound: a self-contained proof}
\label{sec:upper}

We prove the following statement from scratch in this subsection.

\begin{theorem}[Upper bound]
\label{thm:upper}
Let $z>1$ be the solution of $z-3=\log z$, and set
\[
 R=1+\rho,
 \qquad
 \rho=\frac{2}{z-1}=0.5705740966\ldots .
\]
There is a deterministic online algorithm and an absolute constant $C$ such
that every instance with $n$ jobs satisfies
\[
  \ALG\le R\OPT+Cn.
\]
\end{theorem}

The proof has three ingredients.  First, pairwise accounting gives an exact
formula for the cost of a schedule.  Second, a simple convexity argument says
that it is enough to consider four boundary outcomes for every test.  Third,
we give a nonnegative \emph{bank balance} $W$: whenever a boundary outcome
creates positive excess cost, the balance drops by at least that amount, up
to a constant rounding loss.  Thus the proof below uses derivatives only as
shorthand for the first-order change of this bank balance; no optimization
background is needed.

\paragraph{The algorithm.}
For a parameter $c>0$, the formal procedure
\textsc{AdaptiveThreshold}$(c)$ is given in
\cref{alg:adaptive-threshold}.  Immediately before a test, let
\begin{itemize}
\item $x$ be the number of jobs not yet tested, including the job about to be
      tested;
\item $N$ be the number of earlier tests that revealed a strictly positive
      processing time; and
\item $D$ be the number of these $N$ jobs that the algorithm deferred.
\end{itemize}
Define the normalized state
\begin{equation}
\label{eq:upper-y}
 y=\frac{D-RN}{x}.
\end{equation}
Because $D\le N$ and $R>1$, we always have $y\le0$.  For later reuse define
the parameterized threshold
\begin{equation}
\label{eq:upper-A}
 a_c(y)=
 \begin{cases}
  1, & y\le-1,\\[1mm]
  \displaystyle
  1+\frac12\log\!\left(\frac{c+2}{c-2y}\right),
       & -1\le y\le0.
 \end{cases}
\end{equation}
The theorem uses $c=\rho$; throughout this subsection we abbreviate
$A(y)=a_\rho(y)$ and $R=1+\rho$.
The defining equation for $z=(\rho+2)/\rho$ says
\begin{equation}
\label{eq:upper-root-rho}
 \log\frac{\rho+2}{\rho}=\frac2\rho-2.
\end{equation}
Consequently, $A$ is continuous at $-1$, increases from $A(-1)=1$ to
$A(0)=1/\rho=1.7526207478\ldots$, and
\begin{equation}
\label{eq:upper-A-prime}
 A'(y)=\frac{1}{\rho-2y}
 \qquad (-1<y<0).
\end{equation}

\begin{algorithm}[H]
\caption{\textsc{AdaptiveThreshold}$(J,c)$}
\label{alg:adaptive-threshold}
\begin{algorithmic}
\Statex \textbf{Input:} jobs $J$ and parameter $c>0$.
\State Fix an arbitrary test order $j_1,\ldots,j_n$.
\State Set $N\gets0$, $D\gets0$, and $Q\gets\varnothing$.
\State \textbf{For} $i=1,\ldots,n$ \textbf{do}:
\State \quad Set $x\gets n-i+1$ and $y\gets(D-(1+c)N)/x$.
\State \quad Set $a\gets a_c(y)$ according to \cref{eq:upper-A}.
\State \quad Test $j_i$ for one unit and observe $p_{j_i}$.
\State \quad \textbf{If} $p_{j_i}\le a$, process $j_i$ immediately.
\State \quad \textbf{Otherwise}, add $j_i$ to $Q$ and set $D\gets D+1$.
\State \quad \textbf{If} $p_{j_i}>0$, set $N\gets N+1$.
\State Process the jobs of $Q$ in nondecreasing order of their revealed
       processing times.
\end{algorithmic}
\end{algorithm}

For $c=\rho$, the state in the pseudocode is exactly
\eqref{eq:upper-y}.  A genuine zero increments neither counter and completes
at the end of its test.  The procedure never uses raw execution; the same
pseudocode will therefore apply to the large finite-$u$ regimes.

The numerator $D-RN$ in \eqref{eq:upper-y} is best viewed as a balance,
not as an estimate of a processing time.  Every revealed positive job reduces
the balance by $R$; deferring it returns one unit, because every later unit
test will also delay that deferred job.  Dividing by $x$ spreads this balance
over the tests that remain.  A state close to $0$ can afford a threshold
larger than one.  Once the normalized balance falls below $-1$, the proof can
only support the conservative threshold one, so $A$ becomes flat.

We will later derive \eqref{eq:upper-A}, rather than asking the reader to
guess why this particular logarithm appears.

\paragraph{Pairwise accounting.}
For two jobs $i<j$ in test order, charge to their unordered pair the duration
of operations of either job performed while the other job is unfinished.  If
job $i$ is processed immediately, this charge in the algorithm is $1+p_i$.
If $i$ is deferred and $j$ is immediate, it is $2+p_j$.  If both are
deferred, it is $2+\min\{p_i,p_j\}$.  Hence
\begin{equation}
\label{eq:upper-alg-pair}
 \ALG=\sum_i(1+p_i)+\sum_{i<j}g_{ij},
 \qquad
 g_{ij}=
 \begin{cases}
  1+p_i, & i\text{ immediate},\\
  2+p_j, & i\text{ deferred},\ j\text{ immediate},\\
  2+\min\{p_i,p_j\}, & i,j\text{ deferred}.
 \end{cases}
\end{equation}
The offline optimum tests and immediately processes jobs in nondecreasing
order of $p_i$, and therefore
\begin{equation}
\label{eq:upper-opt-pair}
 \OPT=\sum_i(1+p_i)
       +\sum_{i<j}\bigl(1+\min\{p_i,p_j\}\bigr).
\end{equation}
For completeness, the latter claim follows because the $k$th completed job
cannot finish before the machine has performed $k$ tests and the processing
of the $k$ shortest jobs.  The shortest-first schedule attains all these
lower bounds simultaneously.

We bound the excess
\[
 F=\ALG-R\OPT.
\]

\paragraph{Why four boundary outcomes suffice.}
Fix only the symbolic decisions made along a run: for each job, specify
whether its value is zero, positive and immediate, or positive and deferred.
This fixes all later counters and thresholds.  The allowed values of a
coordinate are then, respectively,
\[
 \{0\},\qquad (0,A(y)],\qquad (A(y),\infty).
\]
If all processing times except $p_i$ are fixed inside such a symbolic run,
\eqref{eq:upper-alg-pair}--\eqref{eq:upper-opt-pair} express $F$ as an
affine function of $p_i$ plus negative multiples of
$\min\{p_i,h\}$.  Since $p\mapsto\min\{p,h\}$ is concave, $F$ is convex in
$p_i$.  A convex function on $(0,A(y)]$ has its supremum at one of the two
ends.  On $(A(y),\infty)$, the eventual slope is $1-R=-\rho<0$; convexity
then implies that every earlier slope is also negative, so the supremum is at
the lower end.

It is therefore enough to consider the following four formal outcomes.
\begin{center}
\begin{tabular}{ccll}
\toprule
type & limiting value & algorithmic decision & counter update\\
\midrule
$\mathsf Z$ & $0$ & immediate & none\\
$\mathsf E$ & $0^+$ & immediate & $N\leftarrow N+1$\\
$\mathsf I$ & $A(y)$ & immediate & $N\leftarrow N+1$\\
$\mathsf Q$ & $A(y)^+$ & deferred
             & $N\leftarrow N+1$, $D\leftarrow D+1$\\
\bottomrule
\end{tabular}
\end{center}
Here $0^+$ is not a new processing time.  It means that a positive immediate
value tends to zero \emph{after its symbolic decision has been fixed}.
Although its processing contribution vanishes, it still increments $N$ and
therefore changes all future thresholds.  This is why $0$ and $0^+$ cannot be
identified.  Similarly, $A(y)^+$ approaches the threshold from the deferred
side.  Applying the one-coordinate argument successively shows that a uniform
bound on all formal four-type words implies the same bound on all genuine
instances.  Because $F$ is a finite sum of affine and minimum terms on a fixed
symbolic path, the simultaneous one-sided limit at every formal vertex exists
and is independent of the rates at which its open coordinates approach their
endpoints.

One more elementary property will simplify every minimum in the pairwise
formulas.  Put $u=D-RN\le0$.  After an outcome of type $\mathsf Z$,
$\mathsf E$ or $\mathsf I$, or $\mathsf Q$, respectively, the pair $(u,x)$
changes as follows:
\[
 (u,x-1),\qquad (u-R,x-1),\qquad (u-R,x-1),\qquad
 (u-\rho,x-1).
\]
Whenever $x\ge2$, the new quotient is at most $u/x$.  Thus the states, and
hence the thresholds, seen by successive tests never increase.  When $x=1$
there is no next test and therefore no next threshold to compare.  The
nonzero boundary values $A(y)$ consequently arrive in nonincreasing order.
In particular, the minimum of the current nonzero boundary value and any
earlier one is always the current value.

\paragraph{The proof state and the local excess.}
The algorithm needs only $(x,N,D)$.  The analysis splits $N$ into two kinds.
Among earlier formal outcomes, let
\begin{itemize}
\item $S$ be the number of \emph{substantive} positive outcomes
      $\mathsf I,\mathsf Q$;
\item $e$ be the number of infinitesimal outcomes $\mathsf E$; and
\item $d$ be the number of deferred outcomes $\mathsf Q$.
\end{itemize}
Thus $N=S+e$ and $D=d$.  Define
\begin{equation}
\label{eq:upper-eta-b}
 \eta=\frac{d-RS}{x},
 \qquad
 b=\frac e x,
 \qquad
 y=\eta-Rb.
\end{equation}
We have $d\le S$, hence $\eta\le0$.  The variable $b$ records exactly the
part of the algorithmic state that is invisible to limiting processing-cost
minima: an infinitesimal value affects $N$, but its minimum with every other
processing time tends to zero.

Consider the current boundary outcome at a pre-state $(x,S,e,d)$, and write
$a=A(y)$.  Pairwise accounting gives
\begin{equation}
\label{eq:upper-exact-decomposition}
 F=-\rho\frac{n(n+1)}2+\sum_i\widetilde\ell_i,
\end{equation}
where the exact contribution of the current outcome is
\begin{align}
 \widetilde\ell_{\mathsf Z}
 =\widetilde\ell_{\mathsf E}&=d,\label{eq:upper-exact-ZE}\\
 \widetilde\ell_{\mathsf I}
 &=d+a[x+d-RS-R],\label{eq:upper-exact-I}\\
 \widetilde\ell_{\mathsf Q}
 &=d+a[d-RS-\rho].\label{eq:upper-exact-Q}
\end{align}
Here is a quick derivation.  The current unit test gives one extra unit of
algorithmic pair cost for each earlier deferred job, producing $d$.  Its own
test and its pairs with all earlier jobs give a universal excess
$-\rho(i+1)$ at step $i$; summing this term produces the first term in
\eqref{eq:upper-exact-decomposition}.  For an $\mathsf I$ outcome, the
current processing operation contributes $a(x+d)$ to the algorithm and
$Ra(S+1)$ to the scaled optimum.  For a $\mathsf Q$ outcome, the corresponding
quantities are $a(d+1)$ and $Ra(S+1)$.  Monotonicity of the thresholds justifies
using the current $a$ in every earlier processing-time minimum.

For the potential calculation it is convenient to discard the two favorable
diagonal terms $-Ra$ and $-\rho a$.  Define the relaxed local rewards
\begin{align}
 \ell_{\mathsf Z}=\ell_{\mathsf E}&=d,\label{eq:upper-ell-ZE}\\
 \ell_{\mathsf I}&=d+a(x+d-RS)
                  =d+ax(1+\eta),\label{eq:upper-ell-I}\\
 \ell_{\mathsf Q}&=d+a(d-RS)
                  =d+ax\eta.\label{eq:upper-ell-Q}
\end{align}
Then $\widetilde\ell_c\le\ell_c$ for all four types, and
\begin{equation}
\label{eq:upper-F-relaxed}
 F\le-\rho\frac{n(n+1)}2+\sum_i\ell_i.
\end{equation}

\paragraph{How the logarithmic threshold is found.}
Before giving the full bank balance, it is useful to see where the formula
comes from.  Ignore $0^+$ outcomes for a moment, so $b=0$, and look in the
range $-1\le y\le0$.  Suppose the nonlinear part of the balance has the form
$x^2f(y)$.  Asking the immediate boundary $\mathsf I$ and the deferred
boundary $\mathsf Q$ to be exactly paid by the first-order decrease of this
balance gives
\begin{align}
 -2f+(y-R)f'+A(1+y)&=0,\label{eq:upper-design-I}\\
 1-2f+(y-\rho)f'+Ay&=0.\label{eq:upper-design-Q}
\end{align}
Subtracting the equations gives $f'=A-1$.  Differentiating either equation
then gives
\[
 1+(2y-\rho)A'(y)=0.
\]
With the natural joining condition $A(-1)=1$, this is exactly
\eqref{eq:upper-A}--\eqref{eq:upper-A-prime}.

The universal negative term in \eqref{eq:upper-F-relaxed} has leading
part $-\rho n^2/2$.  The initial balance will be $n^2f(0)$, so exact
cancellation requires $f(0)=\rho/2$.  Since
\begin{equation}
\label{eq:upper-f-definition}
 f(y)=\int_{-1}^{y}(A(t)-1)\,dt,
\end{equation}
direct integration gives
\[
 f(0)=\frac12-\frac{\rho}{4}
       \log\!\left(\frac{\rho+2}{\rho}\right).
\]
Thus $f(0)=\rho/2$ is equivalent to
\eqref{eq:upper-root-rho}, or, after setting $z=(\rho+2)/\rho$, to
$z-3=\log z$.  Thus both the threshold and the constant $R$ are forced by
two local balance equations and the initial accounting; they are not the
output of an unexplained numerical optimization.

\paragraph{The bank balance.}
For $-1\le y\le0$, put
\begin{equation}
\label{eq:upper-fh}
 f(y)=\int_{-1}^{y}(A(t)-1)\,dt,
 \qquad
 h(y)=A(y)(1+y).
\end{equation}
The term $x d$ in the balance below pays for the fact that each of the $x$
remaining tests will delay each of the $d$ already deferred jobs.  The
$x^2f(y)$ term handles substantive outcomes.  Finally, the terms involving
$b$ repair the discontinuity between a genuine zero and $0^+$.

Define
\begin{equation}
\label{eq:upper-W}
 W(x,S,e,d)=x d+x^2G(y,b),
\end{equation}
using two formulas for $G$:
\begin{equation}
\label{eq:upper-G}
 G(y,b)=
 \begin{cases}
  \displaystyle f(y)+b h(y)+\frac R2b^2,
       & y\ge-1,\\[2mm]
  \displaystyle g(\eta),\qquad
       g(\eta)=\frac{(1+\eta)_+^2}{2R},
       & y\le-1,
 \end{cases}
 \qquad \eta=y+Rb.
\end{equation}
The first formula is used while the threshold is changing.  The second is
used once the threshold has saturated at one.  In that flat region the
limiting processing-cost terms depend on the substantive balance $\eta$, not
on the number of $0^+$ outcomes separately, which explains why $g$ contains
no $b$.

This piecewise definition is smooth enough for a first-order argument.  At
$y=-1$, feasibility gives $0\le b\le1/R$ and $\eta=-1+Rb$.  Since
\[
 f(-1)=h(-1)=0,
 \qquad h'(-1)=1,
\]
the first formula in \eqref{eq:upper-G} has value $Rb^2/2$ and derivatives
$G_y=b$, $G_b=Rb$.  The second formula has exactly the same value and
derivatives.  This also checks the derivative in the raw $x$ coordinate:
at fixed $(S,e,d)$,
\[
 \partial_x\!\left[x^2G(y,b)\right]
 =x\left(2G-yG_y-bG_b\right),
\]
which agrees on the two sides because $G$, $G_y$, and $G_b$ agree.  The
remaining raw derivatives are linear combinations of $G_y$ and $G_b$.  The
internal join of $(1+\eta)_+^2$ at $\eta=-1$ also has a continuous first
derivative.  Hence $W$ is continuously differentiable.
Moreover, $W\ge0$: in the first region $f\ge0$, $h\ge0$, and $b,d\ge0$; in
the second this is immediate from \eqref{eq:upper-G}.

The choice $h=A(1+y)$ also has a direct interpretation.  At $b=0$, it makes
the first infinitesimal outcome $0^+$ exactly neutral in the local balance.
The quadratic correction $Rb^2/2$ then maintains the balance for all $b$ and
is precisely what makes the two regions meet smoothly.

\paragraph{Four local balance checks.}
We now verify that every formal outcome can be paid from the bank.  In raw
coordinate order $(x,S,e,d)$, the four unit-step directions are
\begin{equation}
\label{eq:upper-directions}
 v_{\mathsf Z}=(-1,0,0,0),\quad
 v_{\mathsf E}=(-1,0,1,0),\quad
 v_{\mathsf I}=(-1,1,0,0),\quad
 v_{\mathsf Q}=(-1,1,0,1).
\end{equation}
For an endpoint type $q\in\{\mathsf Z,\mathsf E,\mathsf I,\mathsf Q\}$,
define its first-order balance
\begin{equation}
\label{eq:upper-drift}
 J_q=\ell_q+\nabla W\cdot v_q.
\end{equation}
The inequality $J_q\le0$ simply says: local reward plus first-order change in
the bank balance is nonpositive.

We need only a few one-variable identities.  On $-1\le y\le0$,
\begin{align}
 f'&=A-1, & f(-1)&=0, & f(0)&=\frac\rho2,
                                                        \label{eq:upper-canonical-1}\\
 -2f+(y-R)f'+A(1+y)&=0,                 \label{eq:upper-canonical-I}\\
 1-2f+(y-\rho)f'+Ay&=0,                 \label{eq:upper-canonical-Q}\\
 -2f+yf'&=A(\rho-y)-R\le0,              \label{eq:upper-canonical-Z}\\
 h+(\rho-y)h'&=AR+(\rho-y)(1+y)A'\ge AR.
                                                        \label{eq:upper-canonical-h}
\end{align}
The first line follows from the definition and
\eqref{eq:upper-root-rho}.  For
\eqref{eq:upper-canonical-I}--\eqref{eq:upper-canonical-Q}, differentiate
the left-hand sides and use $A'=1/(\rho-2y)$; both derivatives vanish, and
both expressions equal zero at $y=-1$.  Equation
\eqref{eq:upper-canonical-Z} follows from
\eqref{eq:upper-canonical-I}.  Its right-hand side is zero at $-1$ and
decreases afterward, because
\[
 \frac{d}{dy}\bigl[A(y)(\rho-y)\bigr]
 =\frac{\rho-y}{\rho-2y}-A(y)\le0.
\]
Indeed, $y\le0$ gives $(\rho-y)/(\rho-2y)\le1$, while
$A(y)\ge1$.
Finally, \eqref{eq:upper-canonical-h} is obtained by differentiating
$h=A(1+y)$.  Notice also that $h,h'\ge0$.

In the changing-threshold region $y\ge-1$, direct differentiation of
\eqref{eq:upper-W} first gives, for an arbitrary $G(y,b)$,
\begin{align}
 J_{\mathsf Z}/x
   &=-2G+yG_y+bG_b,\label{eq:upper-drift-general-Z}\\
 J_{\mathsf E}/x
   &=-2G+(y-R)G_y+(b+1)G_b,\label{eq:upper-drift-general-E}\\
 J_{\mathsf I}/x
   &=-2G+(y-R)G_y+bG_b+A(1+y+Rb),
                                      \label{eq:upper-drift-general-I}\\
 J_{\mathsf Q}/x
   &=1-2G+(y-\rho)G_y+bG_b+A(y+Rb).
                                      \label{eq:upper-drift-general-Q}
\end{align}
Substituting $G=f+bh+Rb^2/2$ and using
\eqref{eq:upper-canonical-I}--\eqref{eq:upper-canonical-Q} reduces the
four expressions to
\begin{align}
 J_{\mathsf Z}/x
   &=(-2f+yf')+b(-h+yh'),\label{eq:upper-drift-active-Z}\\
 J_{\mathsf E}/x
   &=b[-h+(y-R)h'+R],\label{eq:upper-drift-active-E}\\
 J_{\mathsf I}/x
   &=b[-h+(y-R)h'+AR],\label{eq:upper-drift-active-I}\\
 J_{\mathsf Q}/x
   &=b[-h+(y-\rho)h'+AR].\label{eq:upper-drift-active-Q}
\end{align}
Now the signs are transparent.  The $\mathsf Z$ line is nonpositive by
\eqref{eq:upper-canonical-Z}, $h,h'\ge0$, and $y\le0$.  The $\mathsf Q$
line is nonpositive by \eqref{eq:upper-canonical-h}.  The
$\mathsf I$ line is no larger because $y-R<y-\rho$ and $h'\ge0$; the
$\mathsf E$ line is no larger again because $A\ge1$.  Thus all four checks
hold in the changing-threshold region.

In the flat region $y\le-1$, we have $A=1$ and $G=g(\eta)$.  The four checks
become
\begin{align}
 J_{\mathsf Z}/x=J_{\mathsf E}/x
   &=-2g+\eta g',\label{eq:upper-drift-flat-ZE}\\
 J_{\mathsf I}/x
   &=-2g+(\eta-R)g'+1+\eta,\label{eq:upper-drift-flat-I}\\
 J_{\mathsf Q}/x
   &=1-2g+(\eta-\rho)g'+\eta.\label{eq:upper-drift-flat-Q}
\end{align}
If $\eta\le-1$, then $g=g'=0$, and these expressions are $0$, $1+\eta$,
and $1+\eta$.  If $-1\le\eta\le0$, then $g'=(1+\eta)/R$, and they are,
respectively,
\[
 -\frac{1+\eta}{R},\qquad
 -\frac{1+\eta}{R},\qquad
 0.
\]
They are therefore nonpositive in both cases.  We have proved
\begin{equation}
\label{eq:upper-all-drift}
 J_{\mathsf Z},J_{\mathsf E},J_{\mathsf I},J_{\mathsf Q}\le0
\end{equation}
at every feasible state.

\paragraph{From infinitesimal checks to unit steps.}
Equation \eqref{eq:upper-all-drift} compares a local reward with the
\emph{derivative} of $W$, whereas one job changes the integer counters by one.
This loses only a constant per job.  More precisely, there is an absolute
constant $C_0$ such that every admissible transition satisfies
\begin{equation}
\label{eq:upper-Taylor}
 W(s+v_q)-W(s)\le \nabla W(s)\cdot v_q+C_0.
\end{equation}
We include the bounded-curvature argument because it is what makes the
$O(n)$ term uniform over all instances.

Suppose first that $x\ge2$ and linearly interpolate the unit transition.  The
numerators of both $y$ and $\eta$ have the form $u-\kappa t$, where $u\le0$
and $\kappa\ge0$, while the denominator is $x-t$.  Their derivatives are
therefore nonpositive, so the segment crosses each joining surface at most
once.  In the changing-threshold region, $y\ge-1$ implies
\[
 R(S+e)-d
   =\rho S+(S-d)+Re\le x.
\]
Consequently,
\[
 0\le \frac Sx\le\frac1\rho,
 \qquad
 0\le \frac ex\le\frac1R,
 \qquad
 0\le \frac dx\le\frac1\rho.
\]
After normalizing by $x$, all reachable states in this region thus lie in a
fixed compact set.  The balance $W$ is homogeneous of degree two, so its
second derivatives in the raw coordinates are homogeneous of degree zero and
are uniformly bounded on this normalized set.

In the flat region the formula is even simpler:
\begin{equation}
\label{eq:upper-W-flat-explicit}
 W=x d+\frac{(x+d-RS)_+^2}{2R}.
\end{equation}
Its gradient is globally Lipschitz.  Because the two regional formulas have
the same first derivative at their interface, an interpolated unit step has a
uniformly bounded second derivative except at finitely many harmless joining
points.  More explicitly, let $\phi(t)=W(s+t v_q)$.  The preceding bounds and
the $C^1$ gluing make $\phi'$ Lipschitz with one uniform constant $L$ along
the whole segment, so
\[
 \left|\phi(1)-\phi(0)-\phi'(0)\right|
 \le \int_0^1\left|\phi'(t)-\phi'(0)\right|\,dt
 \le \int_0^1 Lt\,dt=\frac L2.
\]
This proves \eqref{eq:upper-Taylor} for $x\ge2$.  When $x=1$, the
changing-threshold pre-states lie in the same compact normalized set, while
\eqref{eq:upper-W-flat-explicit} applies globally in the flat region.  At
$x=0$ we define $W=0$; this agrees with
\eqref{eq:upper-W-flat-explicit}, since terminal feasibility gives
$d\le S$ and hence $d-RS\le0$.

Combining \eqref{eq:upper-drift}--\eqref{eq:upper-Taylor} gives the actual
one-step bank inequality
\begin{equation}
\label{eq:upper-one-step}
 \ell_i+W(s_{i+1})-W(s_i)\le C_0.
\end{equation}

\paragraph{Telescoping the account.}
Initially $S=e=d=0$, $x=n$, and $y=0$.  Therefore
\begin{equation}
\label{eq:upper-W-initial}
 W(s_0)=n^2f(0)=\frac\rho2n^2.
\end{equation}
The terminal balance is zero.  Summing
\eqref{eq:upper-one-step} over all $n$ jobs yields
\[
 \sum_i\ell_i\le \frac\rho2n^2+C_0n.
\]
Substitution into \eqref{eq:upper-F-relaxed} gives
\[
 \ALG-R\OPT=F
 \le
 -\rho\frac{n(n+1)}2+\frac\rho2n^2+C_0n
 =\left(C_0-\frac\rho2\right)n.
\]
The constant is independent of the four-type word.  Taking
$C=\max\{0,C_0-\rho/2\}$, the boundary reduction therefore extends the same
estimate to every vector of nonnegative processing times, proving
\cref{thm:upper}.

Finally, $\OPT\ge n(n+1)/2$, so the additive $O(n)$ term is lower order and
\[
 \frac{\ALG}{\OPT}\le R+O(1/n).
\]
This proves the desired size-asymptotic upper bound.  Notice exactly where the
remaining finite-instance loss enters: it is only the uniform Taylor
remainder in \eqref{eq:upper-Taylor} (together with the deliberately
discarded favorable diagonal terms), not any unproved limiting step in the
endpoint reduction.  In particular, this argument proves a linear additive
term; it does not prove one fixed additive constant at the endpoint $R$.
